% Independent C.1 bridge supplied by agent /root/homotopy_levels.
% AI-produced by GPT-6 Astra (OpenAI, ChatGPT Pro mode), 2026-11-19.

\paragraph{C1.B1 (Proved: the directed collage at a minimal stratum).}
Let $P$ be a finite poset and let $\mathcal C_\bullet$ be a complete
Segal object in spaces, conservative over $N(P)$. Let $i\in P$ be
minimal, let $A=\mathcal C_i$ (an $\infty$-groupoid), and let
$\mathcal D=\mathcal C|_{P\setminus\{i\}}$. Let $\mathcal M$ be the
category of arrows of $\mathcal C$ whose source is in $A$ and whose
target is in $\mathcal D$. Morphisms in $\mathcal M$ are commutative
squares, with the source-side arrow necessarily an equivalence.
Write $q:\mathcal M\to A$ and $p:\mathcal M\to\mathcal D$ for the
source and target functors. Then there is a canonical pushout formula
in $\operatorname{Cat}_\infty$
\[
 \mathcal C\simeq
 A\mathbin{\coprod}_{\mathcal M\times\{0\}}
 (\mathcal M\times[1])
 \mathbin{\coprod}_{\mathcal M\times\{1\}}\mathcal D.
 \tag{C1.B1}
\]
The maps from the two endpoints are $q$ and $p$, respectively.
The comparison is induced by the tautological natural transformation
from the source to the target of an arrow.

\begin{proof}
We verify the universal property on mapping spaces into an arbitrary
$\infty$-category $\mathcal K$. A functor from the right side to
$\mathcal K$ is a pair $F_A:A\to\mathcal K$,
$F_D:\mathcal D\to\mathcal K$, together with a natural transformation
\[
 \eta:F_Aq\Longrightarrow F_Dp
       \quad\hbox{in }\operatorname{Fun}(\mathcal M,\mathcal K).
 \tag{C1.B2}
\]
Endpoint identifications in this description are equivalences; hence
the description is a homotopy pullback of spaces of functors and
arrows, not a set of strict diagrams.

We show that these data are exactly a functor $\mathcal C\to\mathcal K$.
Restriction gives $F_A,F_D$, and evaluation on each crossing arrow
gives $\eta$. Naturality of $\eta$ follows by applying the functor to
the corresponding square in $\mathcal C$.
Conversely, use $F_A,F_D$ on the two full subcategories and use
$\eta_m$ on a crossing arrow $m$. There is no arrow from
$\mathcal D$ to $A$. Thus a composable string has at most one
crossing arrow. Precomposition with an arrow in $A$ and
postcomposition with an arrow in $\mathcal D$ are precisely the
source and target actions encoded by the squares in $\mathcal M$;
(C1.B2) makes their images commute with composition.

For clarity, the same argument includes all higher coherence, as can
be checked on the complete Segal nerves. A mixed simplex has a unique
last vertex in $A$. Its part before that vertex is a simplex in the
groupoid $A$, so it transports the crossing arrow by an equivalence.
The remaining string is the corresponding simplex of $\mathcal M$
followed by its target simplex in $\mathcal D$. Specifying its image
is therefore specifying the cone whose edges are the components of
$\eta$ on that simplex. The face maps are exactly restriction,
source transport, and target composition. The definition of a
natural transformation supplies these cones compatibly on every
simplex of $\mathcal M$, including their faces and degeneracies.
The Segal equivalences identify the spaces of such extensions with
these spaces of compatible cones; their inverse choices have
contractible fibres. Thus restriction and the construction just
given are inverse equivalences of mapping spaces, naturally in
$\mathcal K$. Yoneda proves (C1.B1).
\end{proof}

\paragraph{C1.B2 (Proved: identification of the incidence category).}
Suppose the deployment is a finite graded topological category with
object spaces $\Sigma_j$, arrow spaces $\Sigma_{jk}$ for $j<k$,
and identities only within each object space. Suppose its source
maps $s_{jk}:\Sigma_{jk}\to\Sigma_j$ are fibrations and the flag
identities and control equations are
\[
 \Sigma_{ijk}=s_{jk}^{-1}(\Sigma_{ij})\subset\Sigma_{jk},
 \qquad s_{ik}=s_{ij}s_{jk}\quad\hbox{on }\Sigma_{ijk}.
 \tag{C1.B3}
\]
Here $\Sigma_{ij}\hookrightarrow\Sigma_j$ is the target inclusion,
and $s_{jk}$ in the second equation is restricted to
$\Sigma_{ijk}\to\Sigma_{ij}$. Define a topological category $M_i$
with objects $\coprod_{j>i}\Sigma_{ij}$, arrows
$\coprod_{i<j<k}\Sigma_{ijk}$, and the induced source, target and
composition maps. Then, after passing to complete Segal objects,
$M_i$ is the $\mathcal M$ of C1.B1. Its target functor
$p:M_i\to\mathcal D$ is a left fibration. Its source functor
$q:M_i\to\Sigma_i$ is well defined by the control equation.

\begin{proof}
The target functor forgets $i$ and regards $\Sigma_{ij}$ as a
subspace of $\Sigma_j$. The flag identities identify the lifting
comparison in each degree with
\[
 (M_i)_n\longrightarrow (M_i)_0
                  \mathbin{\times^h}_{\mathcal D_0}\mathcal D_n.
 \tag{C1.B4}
\]
This is an equivalence. Indeed, the strict fibre product is the
specified iterated flag intersection; it computes the homotopy
fibre product because the source projection of each composable
string is a composite of pullbacks of the given fibrations.
The comparisons in (C1.B4), together with the Segal equivalences,
are precisely the left-fibration lifting condition for simplicial
spaces. Equivalently, an outgoing arrow from $i$ has, for each
further arrow in $\mathcal D$ from its endpoint, a contractible
space of lifts obtained by composition. The source in $\Sigma_i$
stays fixed because of the second identity in (C1.B3).

The same fibre-product calculation identifies the simplices with
those of the category of crossing arrows: in a square between two
such arrows the source-side path is an equivalence in the space
$\Sigma_i$, while the target-side arrows compose in $\mathcal D$.
Replacing strict fibres by homotopy fibres accounts for precisely
these paths. Thus the canonical comparison is an equivalence of
left fibrations over $\mathcal D$, compatible with $q$, and hence
an equivalence of the required crossing-arrow categories.
\end{proof}

In the terminology of \cite[Definition 1.1.4 and Theorem 1.1.7]{B56Haine2023},
the flag diagram is a d\'ecollage: its adjacent-pair Segal comparisons
are equivalences. That theorem identifies such diagrams with
$\infty$-categories conservative over $P$. It confirms the categorical
presentation used here; it does not replace the geometric comparison
with $\operatorname{Exit}(X)$ proved in C1.B3 below.

\paragraph{C1.B2a (Proved: coherent removal of retained outer faces).}
Suppose a finite set $R$ of already removed labels is a lower ideal
in the original poset. Give every retained outer face a product
collar, with all collar coordinates and collar transports commuting
at intersections and with the source maps of the remaining deployment
(whose labels lie outside $R$).
For a deployed space $Z$ write $Z^\circ_R$ for the space with these
outer faces deleted, keeping the remaining stratification.
There is a stratum-preserving homotopy $h^Z_t:Z\to Z$, with
$h^Z_0=\operatorname{id}$ and $h^Z_1(Z)\subset Z^\circ_R$, natural
for the source maps and incoming face inclusions among those remaining
labels. It restricts to a
homotopy of $Z^\circ_R$. Therefore
$Z^\circ_R\hookrightarrow Z$ induces an equivalence on exit categories
when the retained outer faces are ignored as separate strata.

\begin{proof}
Normalize each collar coordinate to $u\in[0,1)$. Choose a continuous
function $f:[0,1)\to(0,1)$ with $f(u)=u$ for $u\geq2/3$ and
$f(0)>0$, and with $f(u)\geq u$. Replace that coordinate at time $t$
by $(1-t)u+tf(u)$ and leave all other coordinates fixed. Extend by
the identity outside the collar. The maps agree at the collar edge.
The finite family of these operations commutes by the product-collar
hypothesis, so their composite is a globally defined homotopy.
It preserves every remaining stratum, since only the ignored outer
coordinates change. At positive time all affected zero coordinates
become positive; at every time originally positive coordinates stay
positive. Compatibility of coordinates and transports makes the
homotopies natural for the stated diagram. The time-one map and
inclusion are thus stratum-preserving homotopy inverses. Applying
exit paths sends such a homotopy to a natural equivalence: each
point's trajectory stays in one stratum, hence is an equivalence,
and the product homotopy supplies its coherence.
\end{proof}

The inclusion of a deleted outer face into its parent is not asserted
to commute with an inward push: its image moves off that face.
Such an incidence is represented by the positive-radius collar map
and its collar homotopy, as in the next paragraph.

For the next minimal label $i$ in $P\setminus R$, let
$q:H_i\to A_i=\Sigma_i$ be its boundary-link bundle. Its old outer
faces are exactly the inverse images of those of $A_i$. Indeed an
$r$-face with $r<i$ pulls back by (C1.B3); an $r$ incomparable with
$i$ has disjoint incoming face; and $i<r$ is impossible because
$R$ is a lower ideal. Hence
\[
 (H_i)^\circ_R=q^{-1}((A_i)^\circ_R).
 \tag{C1.B4a}
\]
In particular a proper bundle $q$ remains proper on these interiors.
The complementary cut space $W$ has the additional outer face $i$.
Its comparison from $(H_i)^\circ_R$ to $W^\circ_{R\cup\{i\}}$
is the positive-radius $i$-collar section, followed by the old
collar pushes. All old pushes commute with this collar, so this
map agrees with the pushed incidence map and the comparison
homotopy is the specified $i$-collar motion. Thus deleting outer
faces preserves the entire gluing diagram coherently, while
preserving properness of the link bundle. We never declare an
ignored boundary to be a separate exit stratum.

\paragraph{C1.B3 (Proved under recursive compatible collar hypotheses).}
Assume, in addition to C1.B2, that the deployment is realized by
proper conically smooth, stratified locally trivial source bundles
and mutually compatible conical collars. Require explicitly that
cutting a minimal stratum produces a collared space $W$
with retained outer boundary $H$, that $H\to\Sigma_i$ is such a
proper bundle in collars, and that both $W$ and $H$ inherit deployments of the
same class with fewer stratum labels. Require that their induced
collars commute with all source maps. This is the recursive
equisingularity condition needed in the induction; mere fibrancy
in C1.B2 is insufficient for this geometric assertion.
Then the complete Segal object of the deployment presents
$\operatorname{Exit}(X)$, naturally with respect to the boundary
inclusions and source maps used in the deployment.

\begin{proof}
Induct on the number of nonempty labels. With one label, the
deployment is its stratum space, presenting its fundamental
$\infty$-groupoid, which is its exit category. Remove a minimal
label $i$. By the boundary-compatible induction hypothesis,
$\mathcal D\simeq\operatorname{Exit}(W)$ and
$M_i\simeq\operatorname{Exit}(H)$, compatibly with
$H\hookrightarrow W$ and $H\to\Sigma_i$.
The punctured cone collar is $H\times(0,\epsilon)$, whose exit
category is $\operatorname{Exit}(H)$ because the interval carries
the trivial stratification and is contractible. Use C1.B2a to remove all previously retained outer faces;
this identifies exit categories and restricts the link bundle by
base change, hence preserves properness. The same collar homotopy
identifies $\operatorname{Exit}(X)$ with the exit category of its
old outer interior. The fibrewise
cone comparison on these interiors gives
\[
 \operatorname{Exit}(C(H\to\Sigma_i))
 \simeq \Sigma_i\mathbin{\coprod}_{M_i\times\{0\}}
                       (M_i\times[1]).
 \tag{C1.B5}
\]
This is exactly \cite[Lemma 2.13]{C_Volpe}, together with
$\operatorname{Exit}(H\times[0,\infty))\simeq
M_i\times[1]$; the boundary point is the smaller stratum.
The geometric cone-collar gluing is an open-cover gluing after
slightly enlarging the collar and shrinking the complementary
cut piece. We first apply C1.B2a: delete the old outer faces
in $H$ and $\Sigma_i$, and all outer faces in $W$. The source
bundle remains proper by (C1.B4a), so Volpe is applied only to
these ordinary conically smooth interiors. The map to the
complementary open piece is the positive-radius collar section
specified in C1.B2a. Its coherent collar homotopies identify this
diagram with the displayed diagram on $W,H,\Sigma_i$. Exit excision,
equivalently \cite[Corollary 2.17]{C_Volpe} or
\cite[Theorem A.7.1]{LurieHA}, therefore gives
\[
 \operatorname{Exit}(X)\simeq
 \operatorname{Exit}(W)
       \mathbin{\coprod}_{\operatorname{Exit}(H)}
 \operatorname{Exit}(C(H\to\Sigma_i)).
 \tag{C1.B6}
\]
Substitute the two induction comparisons and (C1.B5), and compare
with the identical universal pushout (C1.B1). This yields the
equivalence of full exit categories. The construction uses the
source and inclusion maps and their tautological transformations,
so the comparisons commute with the retained boundaries at the
next induction step. The naturality is part of the pushout
universal properties, not an additional choice of independent
equivalences of individual spaces. This completes the induction.
\end{proof}
